
Semantic entropy (SE) provides effective hallucination detection in large language models but requires computationally expensive multi-sample generation. Semantic Entropy Probes (SEPs) have been proposed as efficient single-pass estimators that extract uncertainty information from hidden states, with published results reporting AUROC approximately 0.85 on TruthfulQA. This work attempted to extend the SEP approach with similarity-augmented training, hypothesizing that semantic similarity features would enable cross-model transfer. The experiment yielded negative results: the proposed SE-Distilled Probe (SEDP) achieved Spearman correlation rho = 0.0843 with true semantic entropy (p = 0.283), statistically indistinguishable from zero, and AUROC = 0.52, indistinguishable from random classification. The baseline SEP achieved comparable performance (rho = 0.0835, AUROC = 0.5214). These results represent a 39\% AUROC gap relative to published benchmarks. The experiment tested a single configuration: layer 25 of Llama-3-8B-Instruct, Token-Before-Generation position, and logistic regression probe architecture. The failure may reflect configuration sensitivity rather than fundamental limitations of the probing approach. This work documents the negative result to highlight that SE probe deployment may require systematic validation across layers, token positions, and architectures.
